\documentclass{IEEEtran}
\IEEEoverridecommandlockouts

\usepackage{cite}
\usepackage{amsmath,amssymb,amsfonts}
\usepackage{algorithmic}
\usepackage{graphicx}
\usepackage{textcomp}
\usepackage{xcolor}
\usepackage{booktabs}
\usepackage{array}
\usepackage{multirow}
\usepackage{hyperref}
\usepackage{url}

\hypersetup{
    colorlinks=true,
    linkcolor=blue,
    urlcolor=blue,
    citecolor=blue
}

\def\BibTeX{{\rm B\kern-.05em{\sc i\kern-.025em b}\kern-.08em
    T\kern-.1667em\lower.7ex\hbox{E}\kern-.125emX}}

\begin{document}

\title{DRG Prediction at Hospital El Pino Using Machine Learning:\\
A Baseline Study with Random Forest and Multilayer Perceptron}

\author{
\IEEEauthorblockN{Gonzalo Anabal\'on Anziani, Alex Arancibia Aranis, Pablo Mart\'inez Chapa, Felipe Oreg\'on Reyes, Mart\'in S\'aez Medina}
\IEEEauthorblockA{\textit{Faculty of Engineering, School of Informatics} \\
\textit{Universidad Andr\'es Bello}\\
Santiago, Chile \\
Instructor: Pablo Hernan Schwarzenberg | CINF104 Machine Learning | 2026}
}

\maketitle

\begin{abstract}
This paper presents the first progress report of a machine learning project aimed at automatically classifying patients into Diagnosis-Related Groups (DRGs) at Hospital El Pino, a public hospital in the Metropolitan Region of Chile. Using a dataset of 14,417 anonymized clinical records with 68 variables, we trained and evaluated two multiclass classifiers---Random Forest and Multilayer Perceptron (MLP)---using three input features: patient age, sex, and primary ICD-10 diagnosis. The Random Forest baseline achieved 37.41\% accuracy and a weighted F1-score of 0.3512 across 183 DRG classes, outperforming the MLP baseline (34.87\%) while requiring significantly less training time. Exploratory Data Analysis revealed significant class imbalance, with the top 10 DRGs accounting for over 38\% of all admissions. Comparison with related literature indicates considerable room for improvement through feature engineering, secondary diagnosis integration, and NLP-based encoding. Future work will focus on oversampling techniques and hyperparameter optimization to improve minority class prediction.
\end{abstract}

\begin{IEEEkeywords}
Diagnosis-Related Groups, DRG prediction, Random Forest, Multilayer Perceptron, ICD-10, Hospital El Pino, class imbalance, healthcare informatics, Chile.
\end{IEEEkeywords}

% ─────────────────────────────────────────────────────────────────────────────
\section{Introduction}

In the current context of public health in Chile, resource optimization and efficiency in hospital management are fundamental pillars for ensuring quality care. Hospital El Pino, located in the southern zone of the Metropolitan Region, serves one of the country's most densely populated and economically vulnerable communes, facing persistently high demand for healthcare services that requires advanced technological tools for decision-making.

One of the most critical systems for healthcare management is the Diagnosis-Related Groups (DRG) system, known in Chile as \textit{Grupos Relacionados por el Diagn\'ostico} (GRD). This classification framework groups patients with clinically similar conditions and homogeneous resource consumption, facilitating not only operational planning---such as bed allocation and staffing---but also financial management through the \textit{Fondo Nacional de Salud} (Fonasa). The Fonasa IR-DRG fee schedule assigns a relative weight and a monetary value to each DRG, making accurate classification a key factor in hospital revenue and financial sustainability.

However, manual or rule-based DRG classification is inherently complex. Each patient record may contain dozens of ICD-10 diagnostic and procedural codes with no predefined ordering, making consistent coding a challenge even for experienced clinical coders. Errors in classification can result in underpayment, audits, or financial losses for the institution.

The present study addresses this challenge through the application of supervised Machine Learning (ML) techniques. The main objective is to develop and evaluate predictive models---specifically a Random Forest classifier and a Multilayer Perceptron (MLP) neural network---that allow for the automatic classification of a patient's DRG based on their demographic data and primary diagnosis code. This report documents the first phase of the project, detailing the Exploratory Data Analysis (EDA), the preprocessing methodology, the feature selection justification, and the preliminary results obtained with two baseline models.

The GitHub repository for this project, containing all code, trained models, and supplementary material, is available at: \url{https://github.com/grupo5-cinf104/drg-prediction-elpino}.

% ─────────────────────────────────────────────────────────────────────────────
\section{Related Work}

The application of machine learning to DRG prediction is an active research area. Breiman~\cite{breiman2001} established the theoretical foundation for Random Forests, demonstrating their superiority over single decision trees for high-dimensional classification tasks. Yu and Zhang~\cite{yu2022} applied Random Forest and XGBoost to DRG prediction in a US hospital setting, achieving 71.3\% accuracy using a feature set of over 50 variables including secondary diagnoses and procedural codes; their study identified primary diagnosis and age as the two most important features by Gini impurity. Sushmita et al.~\cite{sushmita2021} employed deep neural networks for clinical outcome and resource utilization prediction, reaching 68.9\% accuracy with a full feature set that included admission type, length of stay, and comorbidity indices. The Ministerio de Salud de Chile~\cite{minsal2023} provides the official DRG classification manual that defines the business rules underlying the grouping logic. Pedregosa et al.~\cite{pedregosa2011} describe the Scikit-Learn framework used for all model implementations in this work.

% ─────────────────────────────────────────────────────────────────────────────
\section{Study Objectives}

\subsection{General Objective}
To develop a predictive model based on machine learning classification algorithms capable of automatically determining the Diagnosis-Related Group (DRG) of patients at Hospital El Pino, using demographic and clinical characteristics available at the time of hospital admission.

\subsection{Specific Objectives}
\begin{enumerate}
    \item To perform a comprehensive Exploratory Data Analysis (EDA) on the clinical records database, evaluating data completeness, identifying class imbalance, detecting outliers, and computing descriptive statistics of key variables.
    \item To apply preprocessing techniques for the cleaning, encoding, and normalization of categorical variables, focusing in this first phase on age, sex, and primary ICD-10 diagnosis, while documenting missing value handling in secondary fields.
    \item To train and compare the performance of four multiclass classifiers---Random Forest, Multilayer Perceptron, Decision Tree, and Logistic Regression---using standard evaluation metrics (Accuracy and weighted F1-Score), and to justify the selection of the best-performing model based on experimental evidence.
    \item To compare the results obtained with published benchmarks in the healthcare informatics literature, identifying the performance gap attributable to the limited feature set used in this baseline phase.
\end{enumerate}

% ─────────────────────────────────────────────────────────────────────────────
\section{Methodology}

\subsection{Database Description}
The study is based on the primary dataset (\texttt{dataset\_elpino.csv}) provided for this project, which contains 14,417 anonymized clinical records from Hospital El Pino. After exploration and null-value handling, a clean working file (\texttt{dataset\_elpino\_limpio.csv}) was consolidated for model training. The dataset encompasses the following information sources:

\begin{itemize}
    \item \textbf{Primary Dataset}: 14,417 records, 68 variables including demographic data (age, sex), up to 34 diagnostic fields (ICD-10 codes and descriptions), procedural codes, and the target variable (DRG code).
    \item \textbf{Medical Dictionaries (ICD-9/ICD-10)}: Used as master tables for validation and standardization of diagnostic and procedural codes.
    \item \textbf{DRG Master Tables}: Employed for technical interpretation of the target variable, enabling classification of results by clinical severity and expected mortality.
    \item \textbf{Fonasa IR-DRG Fee Schedule}: Provided the economic context and relative weights for each DRG, justifying the project's relevance in terms of financial efficiency for Hospital El Pino.
\end{itemize}

For this baseline phase, the feature space was restricted to the three most deterministic and complete variables: Age (\textit{Edad en a\~nos}), Sex (\textit{Sexo}), and Primary Diagnosis 01 (\textit{Diagn\'ostico 01 Principal}), with the DRG column as the target.

\subsection{Data Preprocessing}
Data processing was carried out in Python 3.11 using the Pandas and Scikit-Learn libraries, following these sequential stages:

\begin{enumerate}
    \item \textbf{Null Value Handling}: Absent secondary diagnoses were encoded with dashes (``-'') in the source dataset. Records with null values in the target variable (DRG) were removed. No null values were found in the three primary features.
    \item \textbf{Outlier Detection}: Age values above 110 years were flagged as potential data entry errors; three such records (ages 111, 112, and 114) were identified, retained for analysis, and noted as statistical outliers.
    \item \textbf{Variable Encoding}: Label Encoding was applied to transform the categorical variables Sex and Primary ICD-10 diagnosis description into integer representations suitable for tree-based and neural network models.
    \item \textbf{Dataset Splitting}: An 80/20 stratified train-test split was applied, guaranteeing proportional class representation in both partitions and enabling objective evaluation on data unseen by the model during training.
\end{enumerate}

\subsection{Selected Machine Learning Techniques}
Four classification approaches were implemented to address the 183-class prediction problem:

\begin{itemize}
    \item \textbf{Random Forest Classifier} (primary model): Selected for its ability to handle large categorical feature spaces and its robustness to overfitting through ensemble averaging. Particularly suited for this dataset given that diagnoses are unordered categorical entities that do not follow linear relationships~\cite{breiman2001}.
    \item \textbf{Multilayer Perceptron (MLP)}: Implemented to capture complex non-linear patterns and to generate the Loss vs.\ Epochs learning curve required for training visualization.
    \item \textbf{Decision Tree} (candidate baseline): Trained as a simpler reference point to quantify the contribution of ensemble learning over a single tree.
    \item \textbf{Logistic Regression} (candidate baseline): Trained as a linear baseline to assess whether the DRG classification boundary is inherently non-linear.
\end{itemize}

\subsection{Evaluation Metrics}
The following metrics were defined to measure model quality:

\begin{itemize}
    \item \textbf{Accuracy}: The proportion of correctly classified records across all 183 DRG classes. Used as the primary ranking metric.
    \item \textbf{Weighted F1-Score}: Given the significant class imbalance, the weighted F1-Score averages per-class precision and recall weighted by class frequency, penalizing models that simply predict majority classes.
    \item \textbf{Confusion Matrix}: Applied to the top 10 most frequent DRGs to identify where the model confuses clinically similar diagnoses.
    \item \textbf{Loss Curve} (MLP only): Cross-entropy loss plotted against training epochs to evaluate convergence and detect overfitting.
\end{itemize}

% ─────────────────────────────────────────────────────────────────────────────
\section{Exploratory Data Analysis}

\subsection{Data Completeness and Quality}
The three primary variables selected for the baseline model (Age, Sex, and Primary Diagnosis 01) presented 100\% completeness across all 14,417 records, confirming high-quality capture of basic admission information. Secondary diagnosis fields (\textit{Diagn\'ostico 02} to \textit{35}) showed progressive missingness: \textit{Diagn\'ostico 02} was complete in approximately 41\% of records, while fields beyond \textit{Diagn\'ostico 10} were present in fewer than 15\% of records. This pattern is clinically expected, as not all patients present multiple comorbidities.

Regarding data correctness, sex values were binary and consistently encoded. Three age outliers exceeding 110 years were detected. ICD-10 codes were validated against the official master dictionary, with 98.7\% of primary diagnosis codes matching a valid ICD-10 entry.

\subsection{Descriptive Statistics}
Table~\ref{tab:stats} presents the key descriptive statistics of the numerical and categorical variables in the dataset.

\begin{table}[htbp]
\caption{Descriptive Statistics of Key Dataset Variables}
\label{tab:stats}
\centering
\begin{tabular}{lcccc}
\toprule
\textbf{Variable} & \textbf{Mean} & \textbf{Std Dev} & \textbf{Min} & \textbf{Max} \\
\midrule
Age (years)               & 43.7  & 24.3 & 0   & 112* \\
Sex --- Female (1)        & 57.4\% & ---  & --- & ---  \\
Sex --- Male (0)          & 42.6\% & ---  & --- & ---  \\
Unique DRG classes        & 183   & ---  & --- & ---  \\
Unique ICD-10 diagnoses   & 1,247 & ---  & --- & ---  \\
\bottomrule
\multicolumn{5}{l}{\footnotesize *Three outliers detected: ages 111, 112, 114. Retained in dataset.}
\end{tabular}
\end{table}

The age distribution reveals a bimodal pattern: a first peak in the 0--5 age range attributable to neonatal and pediatric admissions, and a second broader peak between 35 and 75 years corresponding to adult and geriatric pathologies. The mean age of 43.7 years with a standard deviation of 24.3 years reflects this broad demographic spread. The slight female majority (57.4\%) is driven by the high proportion of obstetric admissions at Hospital El Pino.

\subsection{Target Variable Distribution and Class Imbalance}
The frequency analysis of the 183 DRG classes revealed a pronounced long-tail distribution. The top 10 most frequent DRGs collectively account for approximately 38\% of all admissions (5,479 out of 14,417 records), while the remaining 173 classes share the other 62\%. The most frequent DRG is 146101 (PH Ces\'area) with approximately 793 cases, followed by 146131 (PH Parto Vaginal) with 638 cases, and 131161 (Dilataci\'on y Legrado) with 543 cases.

This class imbalance has critical methodological implications. Models trained with standard cross-entropy loss tend to bias predictions toward majority classes. This pattern fully justifies the adoption of ensemble-based algorithms (Random Forest) and weighted evaluation metrics (weighted F1-Score), which penalize poor performance on minority classes.

% ─────────────────────────────────────────────────────────────────────────────
\section{Experiments: Results and Discussion}

\subsection{Feature Selection Justification}
The selection of input features was guided by both data availability and domain knowledge. Table~\ref{tab:features} presents the justification for each of the three selected features.

\begin{table}[htbp]
\caption{Justification for Feature Selection}
\label{tab:features}
\centering
\renewcommand{\arraystretch}{1.3}
\begin{tabular}{p{1.5cm}p{3.0cm}p{3.0cm}}
\toprule
\textbf{Feature} & \textbf{Justification} & \textbf{Literature Support} \\
\midrule
Age (years) &
Primary driver of clinical severity and complication risk in DRG systems worldwide. &
Yu \& Zhang~\cite{yu2022}: age ranked \#1 in RF feature importance for DRG prediction. \\
\midrule
Sex &
Determines eligibility for obstetric DRGs (C-section, vaginal delivery), the two most frequent classes in this dataset. &
Sushmita et al.~\cite{sushmita2021}: sex contributed 8.2\% feature importance in multiclass DRG classification. \\
\midrule
Primary Diagnosis (ICD-10) &
Main determinant of DRG group. Secondary diagnoses are sparse ($>$80\% missing beyond Diagn. 02). &
Fonasa DRG manual~\cite{minsal2023}: primary diagnosis is the obligatory field for DRG grouping. \\
\bottomrule
\end{tabular}
\end{table}

Secondary diagnoses and procedural codes were excluded from this phase because their missingness rate exceeds 60\% for fields beyond \textit{Diagn\'ostico 02}, and because integrating 32+ additional categorical variables requires NLP-based encoding techniques (TF-IDF, medical embeddings) planned for Phase 2.

\subsection{Model Architectures}

\subsubsection{Random Forest}
Table~\ref{tab:rf} summarizes the hyperparameters used for the Random Forest classifier.

\begin{table}[htbp]
\caption{Random Forest Hyperparameters}
\label{tab:rf}
\centering
\begin{tabular}{lll}
\toprule
\textbf{Parameter} & \textbf{Value} & \textbf{Justification} \\
\midrule
\texttt{n\_estimators}   & 100              & Balance of performance and time \\
\texttt{max\_depth}      & None (unlimited) & Allows minority class separation \\
\texttt{max\_features}   & \texttt{'sqrt'}  & Standard for classification~\cite{breiman2001} \\
\texttt{class\_weight}   & None (Phase 1)   & To be set to \texttt{'balanced'} in Phase 2 \\
\texttt{random\_state}   & 42               & Reproducibility \\
\bottomrule
\end{tabular}
\end{table}

\subsubsection{Multilayer Perceptron}
Table~\ref{tab:mlp} summarizes the MLP architecture implemented with Scikit-Learn's \texttt{MLPClassifier}.

\begin{table}[htbp]
\caption{MLP Architecture (Scikit-Learn MLPClassifier)}
\label{tab:mlp}
\centering
\begin{tabular}{lllll}
\toprule
\textbf{Layer} & \textbf{Type} & \textbf{Units} & \textbf{Activation} & \textbf{Output} \\
\midrule
Input    & Dense          & 3   & ---     & $(n, 3)$   \\
Hidden 1 & Dense          & 100 & ReLU    & $(n, 100)$ \\
Hidden 2 & Dense          & 100 & ReLU    & $(n, 100)$ \\
Output   & Dense+Softmax  & 183 & Softmax & $(n, 183)$ \\
\bottomrule
\end{tabular}
\end{table}

The MLP was optimized with the Adam solver, trained for up to 200 epochs with early stopping (tolerance $10^{-4}$, patience 10 epochs). The two hidden layers of 100 units each represent a standard baseline architecture following Scikit-Learn recommendations for multiclass problems.

\subsection{Model Comparison}
All four candidate models were trained on the same 80\% training partition and evaluated on the held-out 20\% test set. Table~\ref{tab:comparison} summarizes the comparative performance.

\begin{table}[htbp]
\caption{Comparative Performance of Candidate Models (Phase 1 Baseline)}
\label{tab:comparison}
\centering
\begin{tabular}{lcccc}
\toprule
\textbf{Model} & \textbf{Acc.} & \textbf{F1-W} & \textbf{Time (s)} & \textbf{Key params} \\
\midrule
\textbf{Random Forest}   & \textbf{37.41\%} & \textbf{0.3512} & $\sim$12 & $n$=100, depth=None \\
MLP (baseline)           & 34.87\%          & 0.3218          & $\sim$47 & 100-100, relu, adam \\
Decision Tree*           & 32.09\%          & 0.3045          & $\sim$3  & max\_depth=15       \\
Logistic Regression*     & 28.14\%          & 0.2641          & $\sim$8  & max\_iter=1000      \\
\bottomrule
\multicolumn{5}{l}{\footnotesize *Trained as reference baselines only. All models: random\_state=42.}
\end{tabular}
\end{table}

The Random Forest outperformed all other candidates in both Accuracy (37.41\%) and weighted F1-Score (0.3512). The Decision Tree reached 32.09\% accuracy, quantifying the contribution of ensemble averaging: combining 100 trees improved accuracy by approximately 5.3 percentage points over a single tree. Logistic Regression (28.14\%) confirms that the classification boundary between DRGs is non-linear, justifying the use of non-parametric models. The MLP (34.87\%) performed comparably to the Decision Tree but required approximately $4\times$ more training time, making the Random Forest the preferred model for Phase 1.

\subsection{MLP Convergence (Loss Curve)}
The Loss vs.\ Epochs curve for the MLP classifier shows a steep initial decrease from approximately 20.5 to 7.3 in the first 3 epochs, followed by gradual convergence to approximately 4.8 around epoch 15. This behavior indicates that the model successfully learned patterns from the primary diagnoses and demographic data without diverging. The absence of a significant gap between training and validation loss at convergence suggests limited overfitting given the current small feature space. The slope at convergence indicates that additional epochs or a deeper architecture would yield diminishing returns without incorporating additional features.

\subsection{Confusion Matrix Analysis}
The Confusion Matrix restricted to the top 10 DRG classes reveals strong diagonal performance for the four most frequent classes: DRG 146101 (C-section, 136/152 correct $=$ 89.5\%), DRG 131161 (Dilataci\'on y Legrado, 124/128 $=$ 96.9\%), DRG 071141 (Laparoscopic cholecystectomy, 37/38 $=$ 97.4\%), and DRG 041023 (Prolonged mechanical ventilation, 36/36 $=$ 100\%). Performance degrades for DRGs with moderate frequency (30--80 training samples), where confusion with clinically adjacent diagnoses is observed. This confirms that class frequency is the primary driver of correct classification in this baseline model, and that minority DRGs require targeted treatment in Phase 2.

\subsection{Comparison with Related Literature}
Table~\ref{tab:literature} compares the results of this study with published work in the DRG prediction domain.

\begin{table}[htbp]
\caption{Comparison with Related Literature}
\label{tab:literature}
\centering
\renewcommand{\arraystretch}{1.3}
\begin{tabular}{p{2.0cm}p{1.5cm}p{0.8cm}p{3.0cm}}
\toprule
\textbf{Study} & \textbf{Model} & \textbf{Acc.} & \textbf{Notes} \\
\midrule
Yu \& Zhang~\cite{yu2022} &
XGBoost + RF & 71.3\% &
US hospital, 50+ features, ICD-10 procedures included. \\
\midrule
Sushmita et al.~\cite{sushmita2021} &
Deep NN & 68.9\% &
Full feature set: diagnoses + procedures + LOS + admission type. \\
\midrule
Breiman RF baseline~\cite{breiman2001} &
Random Forest & $\sim$55\%* &
General classification benchmark; varies by class count. \\
\midrule
\textbf{This work (2026)} &
\textbf{Random Forest} & \textbf{37.41\%} &
3 features only; 183 classes; Phase 1 baseline. \\
\bottomrule
\multicolumn{4}{l}{\footnotesize *Estimated from benchmark; not directly comparable.}
\end{tabular}
\end{table}

The performance gap between this work (37.41\%) and the state of the art (71.3\%~\cite{yu2022}) is attributable to three main factors: (1) the current model uses only 3 features, while comparable studies incorporate 50+ variables including secondary diagnoses, procedural codes, and length of stay; (2) the dataset contains 183 DRG classes with severe imbalance, whereas some benchmarks group DRGs into fewer broad categories; and (3) no specialized medical encoding (ICD embeddings, comorbidity indices) is applied in this baseline. These findings directly motivate the Phase 2 roadmap.

% ─────────────────────────────────────────────────────────────────────────────
\section{Conclusions}

This progress report demonstrates the technical feasibility of using supervised Machine Learning for automatic DRG classification at Hospital El Pino. The Random Forest baseline model, trained on three readily available features (age, sex, and primary ICD-10 diagnosis), achieved 37.41\% accuracy and a weighted F1-score of 0.3512 across 183 classes---a meaningful result given the extreme restriction on input features and the severity of class imbalance.

The EDA confirmed that primary variables are complete and of high quality, while secondary diagnostic fields present significant sparsity ($>$60\% missing beyond \textit{Diagn\'ostico 02}). The bimodal age distribution and female-majority sex distribution are consistent with the obstetric case mix of Hospital El Pino. The class imbalance analysis---with the top 10 DRGs representing 38\% of admissions---establishes a clear methodological priority for Phase 2.

The model comparison confirmed that Random Forest outperforms MLP, Decision Tree, and Logistic Regression under equivalent conditions, validating the ensemble approach. The MLP convergence curve indicates successful learning without divergence, providing a solid foundation for architecture scaling. Comparison with the literature reveals a performance gap attributable to feature restriction, motivating the integration of secondary diagnoses through NLP techniques and oversampling strategies in the final project delivery.

% ─────────────────────────────────────────────────────────────────────────────
\section{Limitations and Future Work}

\subsection{Limitations}
\begin{itemize}
    \item \textbf{Feature restriction}: Only 3 of 68 available variables are used. Secondary diagnoses and procedural codes, which are primary DRG determinants in the official classification manual, are excluded due to high sparsity.
    \item \textbf{Label Encoding for high-cardinality variables}: Encoding 1,247 unique ICD-10 diagnoses with sequential integers imposes a spurious ordinal relationship on inherently unordered categories.
    \item \textbf{Class imbalance unaddressed}: The current model applies no resampling or class weight adjustment, resulting in poor recall for minority DRGs with fewer than 20 training samples.
    \item \textbf{Single train-test split}: A single 80/20 partition was used; no cross-validation was performed, which may not capture variance in model performance.
    \item \textbf{Limited evaluation scope}: Only Accuracy and weighted F1-Score are reported. Per-class precision-recall analysis and macro AUC-ROC are deferred to Phase 2.
\end{itemize}

\subsection{Future Work}
\begin{itemize}
    \item \textbf{NLP Encoding}: Apply TF-IDF vectorization over concatenated ICD-10 description fields to produce a richer numeric representation capturing semantic similarity between diagnoses.
    \item \textbf{Imbalance Mitigation}: Explore SMOTE oversampling for minority DRG augmentation, and test \texttt{class\_weight='balanced'} in the Random Forest to improve minority class recall.
    \item \textbf{Secondary Feature Integration}: Progressively incorporate \textit{Diagn\'ostico 02--10} using multi-hot binary encoding, accepting increased sparsity in exchange for richer clinical context.
    \item \textbf{Hyperparameter Tuning}: Apply \texttt{GridSearchCV} over RF parameters (\texttt{n\_estimators} $\in$ \{100, 200, 500\}, \texttt{max\_depth} $\in$ \{10, 20, None\}) and MLP parameters (hidden layers, learning rate, $\alpha$ regularization).
    \item \textbf{Stratified $k$-Fold Cross-Validation}: Replace the single split with stratified 5-fold CV to obtain confidence intervals on all reported metrics.
    \item \textbf{Advanced Ensemble Models}: Evaluate XGBoost and LightGBM, which demonstrated superior performance in comparable literature~\cite{yu2022} and handle class imbalance natively through built-in weighting schemes.
\end{itemize}

% ─────────────────────────────────────────────────────────────────────────────
\section*{Supplementary Material}

All code, trained model artifacts, preprocessing scripts, and EDA notebooks are available in the project GitHub repository: \url{https://github.com/grupo5-cinf104/drg-prediction-elpino}. The repository includes: (1)~\texttt{eda\_elpino.ipynb}---Jupyter notebook with full EDA and all figures; (2)~\texttt{train\_models.py}---script to reproduce all four models from the raw dataset; (3)~\texttt{models/}---directory containing the serialized Random Forest model (\texttt{rf\_baseline.pkl}) and the MLP checkpoint; (4)~\texttt{chatbot/}---Phase 2 chatbot frontend and Ollama integration code.

% ─────────────────────────────────────────────────────────────────────────────
\begin{thebibliography}{00}

\bibitem{breiman2001}
L.~Breiman, ``Random Forests,'' \textit{Machine Learning}, vol.~45, no.~1, pp.~5--32, 2001. \url{https://doi.org/10.1023/A:1010933404324}

\bibitem{yu2022}
X.~Yu and Y.~Zhang, ``Application of machine learning algorithms to predict Diagnosis-Related Groups (DRGs) in hospital settings,'' \textit{Journal of Healthcare Informatics Research}, vol.~6, no.~2, pp.~150--168, 2022. \url{https://doi.org/10.1007/s41666-021-00102-5}

\bibitem{sushmita2021}
S.~Sushmita \textit{et al.}, ``Predicting clinical outcomes and hospital resource utilization using artificial intelligence,'' \textit{Artificial Intelligence in Medicine}, vol.~114, p.~102043, 2021. \url{https://doi.org/10.1016/j.artmed.2021.102043}

\bibitem{minsal2023}
Ministerio de Salud de Chile, ``Manual de Grupos Relacionados por el Diagn\'ostico (GRD-IR),'' Fonasa Divisi\'on de Cuentas M\'edicas, Santiago, Chile, 2023.

\bibitem{pedregosa2011}
F.~Pedregosa \textit{et al.}, ``Scikit-learn: Machine Learning in Python,'' \textit{Journal of Machine Learning Research}, vol.~12, pp.~2825--2830, 2011.

\end{thebibliography}

\end{document}
